% Zone „Das kennst du schon“ – aus bank/trigonometrie/zone.jsonl (zusammenbau v0.1)
\einheitenkopf[zone][Kennst du schon]{Das kennst du schon}
\setcounter{aufgabe}{0}
%% TODO Grundvorstellungs-Aufgabe als erste Hauptnummer der Zone (2.8) fehlt in der Bank

%% TODO Titel: Ich-kann-Satz fehlt in der Bank (Platzhalter: Kettenname) – Nr. 1
%% TODO Anweisung: ein Satz über der ersten Teilaufgabe fehlt in der Bank – Nr. 1
\begin{aufgabe}{Rechtwinkliges Dreieck erkennen, Rechtwinkelmarke (Bogen mit Punkt), Hypotenuse gegenüber dem rechten Winkel, Katheten am rechten Winkel}
\swz{Das Dreieck hat bei $C$ einen rechten Winkel. Welche Seite ist die Hypotenuse?}{\dreieck{(0,0)}{(5,0)}{(1.8,2.4)}{a}{b}{c}{}{}{}}
%% TODO Darstellung für schwach fehlt in der Bank (trigonometrie-zone-f1-v3; Abschnitt „Für schwache Schüler“)
\swz{Im Dreieck $PQR$ liegt der rechte Winkel bei $Q$. Welche Seite ist die Hypotenuse, welche sind die Katheten?}{}
\end{aufgabe}

%% TODO Titel: Ich-kann-Satz fehlt in der Bank (Platzhalter: Kettenname) – Nr. 2
%% TODO Anweisung: ein Satz über der ersten Teilaufgabe fehlt in der Bank – Nr. 2
\begin{aufgabe}{Winkel benennen (α, β, γ, auch β1, β2, ε), Winkel messen und schätzen, spitz oder stumpf}
\begin{teile}
\teil Ein Winkel ist $68^\circ$ groß. Spitz oder stumpf? \\ \kreuz{spitz} \kreuz{stumpf}
\end{teile}
\begin{teile}
\teil Wie groß ist der Winkel $\varepsilon$ ungefähr? \\ \kreuz{etwa $45^\circ$} \\ \kreuz{etwa $90^\circ$} \\ \kreuz{etwa $135^\circ$}
\end{teile}
\end{aufgabe}

%% TODO Titel: Ich-kann-Satz fehlt in der Bank (Platzhalter: Kettenname) – Nr. 3
%% TODO Anweisung: ein Satz über der ersten Teilaufgabe fehlt in der Bank – Nr. 3
\begin{aufgabe}{Taschenrechner}
%% TODO Darstellung für schwach fehlt in der Bank (trigonometrie-zone-f3-v1; Abschnitt „Für schwache Schüler“)
\swz{Tippe mit Klammern: $\frac{18}{4 + 2}$ – was kommt heraus?}{}
%% TODO Darstellung für schwach fehlt in der Bank (trigonometrie-zone-f3-v3; Abschnitt „Für schwache Schüler“)
\swz{Berechne mit dem Taschenrechner: $\frac{8{,}4}{2{,}1 + 1{,}4}$}{}
\end{aufgabe}

%% TODO Titel: Ich-kann-Satz fehlt in der Bank (Platzhalter: Kettenname) – Nr. 4
%% TODO Anweisung: ein Satz über der ersten Teilaufgabe fehlt in der Bank – Nr. 4
\begin{aufgabe}{Gleichung mit einem Bruch nach einer Größe umstellen, Größe im Zähler (mal) oder im Nenner (geteilt), mit Strich}
%% TODO Darstellung für schwach fehlt in der Bank (trigonometrie-zone-f4-v1; Abschnitt „Für schwache Schüler“)
\swz{$\text{$\frac{x}{5} = 3$ – wie groß ist $x$?}$}{}
%% TODO Darstellung für schwach fehlt in der Bank (trigonometrie-zone-f4-v3; Abschnitt „Für schwache Schüler“)
\swz{$\text{$\frac{x}{2{,}5} = 0{,}6$ – wie groß ist $x$?}$}{}
\end{aufgabe}

%% TODO Titel: Ich-kann-Satz fehlt in der Bank (Platzhalter: Kettenname) – Nr. 5
%% TODO Anweisung: ein Satz über der ersten Teilaufgabe fehlt in der Bank – Nr. 5
\begin{aufgabe}{Bruch als Verhältnis lesen und in eine Dezimalzahl umwandeln}
%% TODO Darstellung für schwach fehlt in der Bank (trigonometrie-zone-f5-v1; Abschnitt „Für schwache Schüler“)
\swz{$\frac{3}{4}$ als Dezimalzahl?}{}
%% TODO Darstellung für schwach fehlt in der Bank (trigonometrie-zone-f5-v3; Abschnitt „Für schwache Schüler“)
\swz{$\frac{7}{8}$ als Dezimalzahl?}{}
\end{aufgabe}

%% TODO Titel: Ich-kann-Satz fehlt in der Bank (Platzhalter: Kettenname) – Nr. 6
%% TODO Anweisung: ein Satz über der ersten Teilaufgabe fehlt in der Bank – Nr. 6
\begin{aufgabe}{Satz des Pythagoras als zweiter Weg zur dritten Seite, Vergleich der Ergebnisse}
%% TODO Darstellung für schwach fehlt in der Bank (trigonometrie-zone-f6-v1; Abschnitt „Für schwache Schüler“)
\swz{Katheten $3$ cm und $4$ cm – wie lang ist die Hypotenuse?}{}
%% TODO Darstellung für schwach fehlt in der Bank (trigonometrie-zone-f6-v3; Abschnitt „Für schwache Schüler“)
\swz{Katheten $2{,}4$ m und $3{,}2$ m – wie lang ist die Hypotenuse?}{}
\end{aufgabe}

%% TODO Titel: Ich-kann-Satz fehlt in der Bank (Platzhalter: Kettenname) – Nr. 7
%% TODO Anweisung: ein Satz über der ersten Teilaufgabe fehlt in der Bank – Nr. 7
\begin{aufgabe}{Winkelsumme im Dreieck, die beiden spitzen Winkel des rechtwinkligen Dreiecks als Ergänzung, Winkel aus Teilwinkeln (Differenz), Strecke aus Teilstrecken}
%% TODO Darstellung für schwach fehlt in der Bank (trigonometrie-zone-f7-v1; Abschnitt „Für schwache Schüler“)
\swz{Ein Dreieck hat die Winkel $45^\circ$ und $85^\circ$. Wie groß ist der dritte?}{}
%% TODO Darstellung für schwach fehlt in der Bank (trigonometrie-zone-f7-v3; Abschnitt „Für schwache Schüler“)
\swz{Der Winkel bei $A$ ist $83^\circ$ groß. Die Strecke $AD$ teilt ihn; der eine Teil ist $47^\circ$. Wie groß ist der andere Teil?}{}
\end{aufgabe}

%% TODO Titel: Ich-kann-Satz fehlt in der Bank (Platzhalter: Kettenname) – Nr. 8
%% TODO Anweisung: ein Satz über der ersten Teilaufgabe fehlt in der Bank – Nr. 8
\begin{aufgabe}{Längeneinheiten vor dem Rechnen angleichen (m und km, cm und m), Ergebnis mit Einheit, Runden auf eine Dezimale}
%% TODO Darstellung für schwach fehlt in der Bank (trigonometrie-zone-f8-v1; Abschnitt „Für schwache Schüler“)
\swz{$3$ km – wie viel Meter?}{}
%% TODO Darstellung für schwach fehlt in der Bank (trigonometrie-zone-f8-v3; Abschnitt „Für schwache Schüler“)
\swz{$1{,}8$ km und $470$ m zusammen – wie viel Kilometer, auf eine Dezimale?}{}
\end{aufgabe}

%% TODO Titel: Ich-kann-Satz fehlt in der Bank (Platzhalter: Kettenname) – Nr. 9
%% TODO Anweisung: ein Satz über der ersten Teilaufgabe fehlt in der Bank – Nr. 9
\begin{aufgabe}{Höhe, Diagonale, Überstand am Trapez, Fußpunkt auf der Verlängerung beim Parallelogramm, halbe Diagonale im Drachen}
%% TODO Darstellung für schwach fehlt in der Bank (trigonometrie-zone-f9-v1; Abschnitt „Für schwache Schüler“)
\swz{Gleichschenkliges Trapez: untere Seite $10$ cm, obere Seite $6$ cm. Wie lang ist ein Überstand?}{}
%% TODO Darstellung für schwach fehlt in der Bank (trigonometrie-zone-f9-v3; Abschnitt „Für schwache Schüler“)
\swz{Gleichschenkliges Trapez: untere Seite $9{,}4$ m, obere Seite $5{,}2$ m. Wie lang ist ein Überstand?}{}
\end{aufgabe}

%% TODO Titel: Ich-kann-Satz fehlt in der Bank (Platzhalter: Kettenname) – Nr. 10
%% TODO Anweisung: ein Satz über der ersten Teilaufgabe fehlt in der Bank – Nr. 10
\begin{aufgabe}{Gleichung mit einem Bruch nach einer Größe umstellen, Größe im Zähler (mal) oder im Nenner (geteilt), mit Strich – Fehler finden}
\swz[3]{Wo ist der Fehler? Rechne richtig. \rechnung{\frac{8}{x} &= 0{,}4 \\ x &= 8 \cdot 0{,}4 \\ x &= 3{,}2}}{}
\end{aufgabe}

%% TODO Titel: Ich-kann-Satz fehlt in der Bank (Platzhalter: Kettenname) – Nr. 11
%% TODO Anweisung: ein Satz über der ersten Teilaufgabe fehlt in der Bank – Nr. 11
\begin{aufgabe}{Gleichung mit einem Bruch nach einer Größe umstellen, Größe im Zähler (mal) oder im Nenner (geteilt), mit Strich – selbst rechnen}
%% TODO Darstellung für schwach fehlt in der Bank (trigonometrie-zone-f4-v6; Abschnitt „Für schwache Schüler“)
\swz{$\text{$\frac{9}{x} = 0{,}6$ – wie groß ist $x$?}$}{}
\end{aufgabe}
